En este trabajo se ha realizado el desarrollo de un IDS basado en aprendizaje automático. A partir del material del conjunto de datos CIC-IDS2017, se ha procesado el tráfico contenido en este y generado características adicionales que describen el payload de los flujos. Estas características son medidas de entropía y proporciones de bytes imprimibles y no ASCII. Tras esto, se ha entrenado un clasificador Random Forest con el objetivo de distinguir entre tráfico benigno y malicioso. Es necesario dividir los resultados en dos experimentos, ya que al analizar el conjunto de datos final, se vio que había muestras con identificadores de flujos repetidos, por lo que al calcular y asignar las nuevas caractertísticas, aquellas muestras que tuvieran el mismo identificador tendrían el mismo valor de características.

Los resultados muestran que para ambos experimentos se obtienen valores elevados de precisión global, aunque ambos pecan de que esto es debido al desbalance entre la clases. Las clases minoritarias tienen más problemas para ser clasificadas correctamente. En general, se obtienen mejores resultados en el experimento 1 aunque tiene el problema del \textit{data leakage}, por lo que el experimento 2 sería más correcto de cara a implantarlo en una red real. Además, a partir del análisis de la importancia de características se ha visto que las que se han añadido han sido de mucha relevancia a la hora de clasificar muestras, superando a todas las que venían por defecto en el conjunto de datos. 

En general, se puede afirmar que se han cumplido los objetivos propuestos en la sección \ref{SEC:OBJETIVOS}. Se ha desarrollado con éxito un sistema de detección de intrusiones basado en aprendizaje automático, más concretamente Random Forest. No obstante, tiene sus limitaciones, las cuales se podrían pulir si se siguiera desarrollando este proyecto. En la siguiente sección se describirá como se podría seguir mejorando el rendimiento de este IDS.